%% toc: P2 11 11.1,11.2,11.3
\section{Proofs for the propositional equivalence laws}
\label{sec:annex-p02-w02-p02-007-equivalence-laws}

Throughout, \(v\) is an arbitrary valuation and \(\widehat v\) its recursive
extension from Section~\ref{sec:prop-syntax-semantics-02}; \(v\entails\alpha\)
means \(\widehat v(\alpha)=\mathbf T\), and \(v\not\entails\alpha\) means
\(\widehat v(\alpha)=\mathbf F\).  Three facts are used repeatedly.

\emph{(F1) Two-valuedness.}  Since \(\widehat v\) takes values in the
two-element set \(\{\mathbf F,\mathbf T\}\), for formulas \(\alpha,\beta\) we
have \(\widehat v(\alpha)=\widehat v(\beta)\) iff
(\(v\entails\alpha\) iff \(v\entails\beta\)), and also iff
(\(v\not\entails\alpha\) iff \(v\not\entails\beta\)).  Hence
\(\alpha\equiv\beta\) holds as soon as one of these biconditionals is shown
for an arbitrary \(v\).

\emph{(F2) The clauses in the form used below.}  \(v\not\entails\bottom\);
\(v\entails\neg\alpha\) iff \(v\not\entails\alpha\);
\(v\entails\alpha\land\beta\) iff \(v\entails\alpha\) and
\(v\entails\beta\); \(v\entails\alpha\lor\beta\) iff \(v\entails\alpha\) or
\(v\entails\beta\); and \(v\not\entails\alpha\to\beta\) iff
\(v\entails\alpha\) and \(v\not\entails\beta\).

\emph{(F3) Verum.}  \(\topv\) is \(\neg\bottom\).  Since
\(v\not\entails\bottom\), the negation clause gives \(v\entails\topv\) for
every \(v\).

\begin{annexitem}{annex:p02-w02-p02-007-dm}{thm:p02-w02-p02-007-de-morgan}{De Morgan laws}
\begin{proof}
\emph{De Morgan for \(\land\).}  Split on the values of \(\varphi,\psi\)
under \(v\).  If \(v\entails\varphi\) and \(v\entails\psi\), then
\(v\entails\varphi\land\psi\), so \(v\not\entails\neg(\varphi\land\psi)\);
also \(v\not\entails\neg\varphi\) and \(v\not\entails\neg\psi\), so
\(v\not\entails\neg\varphi\lor\neg\psi\).  Otherwise
\(v\not\entails\varphi\) or \(v\not\entails\psi\).  Then
\(v\not\entails\varphi\land\psi\), so \(v\entails\neg(\varphi\land\psi)\);
and \(v\entails\neg\varphi\) or \(v\entails\neg\psi\), so
\(v\entails\neg\varphi\lor\neg\psi\).  In both cases the two sides have the
same value, and (F1) gives the law.

\emph{De Morgan for \(\lor\).}  If \(v\not\entails\varphi\) and
\(v\not\entails\psi\), then \(v\not\entails\varphi\lor\psi\), so
\(v\entails\neg(\varphi\lor\psi)\); also \(v\entails\neg\varphi\) and
\(v\entails\neg\psi\), so \(v\entails\neg\varphi\land\neg\psi\).
Otherwise \(v\entails\varphi\) or \(v\entails\psi\).  Then
\(v\entails\varphi\lor\psi\), so \(v\not\entails\neg(\varphi\lor\psi)\);
and \(v\not\entails\neg\varphi\) or \(v\not\entails\neg\psi\), so
\(v\not\entails\neg\varphi\land\neg\psi\).  Again the values agree.
\end{proof}
\end{annexitem}

\begin{annexitem}{annex:p02-w02-p02-007-dist}{thm:p02-w02-p02-007-distribution}{Distribution laws}
\begin{proof}
In each of the four laws split on whether \(v\entails\varphi\).

\emph{\(\land\) over \(\lor\), left.}  If \(v\not\entails\varphi\), then
\(v\not\entails\varphi\land(\psi\lor\chi)\), and
\(v\not\entails\varphi\land\psi\), \(v\not\entails\varphi\land\chi\), so
\(v\not\entails(\varphi\land\psi)\lor(\varphi\land\chi)\).  If
\(v\entails\varphi\), then \(v\entails\varphi\land(\psi\lor\chi)\) iff
\(v\entails\psi\lor\chi\) iff \(v\entails\psi\) or \(v\entails\chi\); and
\(v\entails\varphi\land\psi\) iff \(v\entails\psi\),
\(v\entails\varphi\land\chi\) iff \(v\entails\chi\), so
\(v\entails(\varphi\land\psi)\lor(\varphi\land\chi)\) iff
\(v\entails\psi\) or \(v\entails\chi\).  The two sides agree in both cases.

\emph{\(\land\) over \(\lor\), right.}  The same two cases apply with each
conjunction written in the other order: if \(v\not\entails\varphi\), both
\((\psi\lor\chi)\land\varphi\) and \((\psi\land\varphi)\lor(\chi\land\varphi)\)
are false; if \(v\entails\varphi\), each side is true iff
\(v\entails\psi\) or \(v\entails\chi\), because the conjunction clause of
(F2) is symmetric in its two conjuncts.

\emph{\(\lor\) over \(\land\), left.}  If \(v\entails\varphi\), then
\(v\entails\varphi\lor(\psi\land\chi)\), and \(v\entails\varphi\lor\psi\),
\(v\entails\varphi\lor\chi\), so
\(v\entails(\varphi\lor\psi)\land(\varphi\lor\chi)\).  If
\(v\not\entails\varphi\), then \(v\entails\varphi\lor(\psi\land\chi)\) iff
\(v\entails\psi\land\chi\) iff \(v\entails\psi\) and \(v\entails\chi\);
and \(v\entails\varphi\lor\psi\) iff \(v\entails\psi\),
\(v\entails\varphi\lor\chi\) iff \(v\entails\chi\), so
\(v\entails(\varphi\lor\psi)\land(\varphi\lor\chi)\) iff
\(v\entails\psi\) and \(v\entails\chi\).

\emph{\(\lor\) over \(\land\), right.}  As in the previous case, with each
disjunction written in the other order; the disjunction clause of (F2) is
symmetric in its two disjuncts.  By (F1) all four laws hold.
\end{proof}
\end{annexitem}

\begin{annexitem}{annex:p02-w02-p02-007-abs}{thm:p02-w02-p02-007-absorption}{Absorption laws}
\begin{proof}
If \(v\entails\varphi\), then \(v\entails\varphi\lor\psi\), hence
\(v\entails\varphi\land(\varphi\lor\psi)\); and
\(v\entails\varphi\lor(\varphi\land\psi)\) because its left disjunct holds.
Both left-hand sides are true, as is \(\varphi\).  If
\(v\not\entails\varphi\), then \(v\not\entails\varphi\land(\varphi\lor\psi)\)
because its left conjunct fails; and \(v\not\entails\varphi\land\psi\), so
both disjuncts of \(\varphi\lor(\varphi\land\psi)\) fail and it is false.
Both left-hand sides are false, as is \(\varphi\).  By (F1) both laws hold.
\end{proof}
\end{annexitem}

\begin{annexitem}{annex:p02-w02-p02-007-dn}{thm:p02-w02-p02-007-double-negation}{Double-negation law}
\begin{proof}
Applying the negation clause twice, \(v\entails\neg\neg\varphi\) iff
\(v\not\entails\neg\varphi\) iff it is not the case that
\(v\not\entails\varphi\).  By two-valuedness the last condition is
\(v\entails\varphi\).  By (F1), \(\neg\neg\varphi\equiv\varphi\).
\end{proof}
\end{annexitem}

\begin{annexitem}{annex:p02-w02-p02-007-mi}{thm:p02-w02-p02-007-material-implication}{Material-implication law}
\begin{proof}
Compare the conditions for falsity.  By (F2),
\(v\not\entails\varphi\to\psi\) iff \(v\entails\varphi\) and
\(v\not\entails\psi\).  By the disjunction clause,
\(v\not\entails\neg\varphi\lor\psi\) iff \(v\not\entails\neg\varphi\) and
\(v\not\entails\psi\), and by the negation clause \(v\not\entails\neg\varphi\)
iff \(v\entails\varphi\).  The two falsity conditions coincide, so by (F1)
\(\varphi\to\psi\equiv\neg\varphi\lor\psi\).
\end{proof}
\end{annexitem}

\begin{annexitem}{annex:p02-w02-p02-007-cp}{thm:p02-w02-p02-007-contraposition}{Contraposition law}
\begin{proof}
By (F2), \(v\not\entails\neg\psi\to\neg\varphi\) iff
\(v\entails\neg\psi\) and \(v\not\entails\neg\varphi\).  By the negation
clause this is: \(v\not\entails\psi\) and \(v\entails\varphi\).  This is
exactly the condition for \(v\not\entails\varphi\to\psi\) given by (F2).  The
falsity conditions coincide, so by (F1)
\(\varphi\to\psi\equiv\neg\psi\to\neg\varphi\).
\end{proof}
\end{annexitem}

\begin{annexitem}{annex:p02-w02-p02-007-lat}{thm:p02-w02-p02-007-lattice}{Commutativity, associativity and idempotence}
\begin{proof}
\emph{Commutativity.}  \(v\entails\varphi\land\psi\) iff
\(v\entails\varphi\) and \(v\entails\psi\) iff \(v\entails\psi\land\varphi\);
likewise \(v\entails\varphi\lor\psi\) iff one of \(\varphi,\psi\) holds under
\(v\) iff \(v\entails\psi\lor\varphi\).

\emph{Associativity.}  Applying the conjunction clause twice,
\(v\entails(\varphi\land\psi)\land\chi\) iff all three of
\(\varphi,\psi,\chi\) hold under \(v\), and the same condition results for
\(\varphi\land(\psi\land\chi)\).  Applying the disjunction clause twice,
\(v\not\entails(\varphi\lor\psi)\lor\chi\) iff none of the three holds under
\(v\), and the same condition results for \(\varphi\lor(\psi\lor\chi)\).

\emph{Idempotence.}  \(v\entails\varphi\land\varphi\) iff
\(v\entails\varphi\) and \(v\entails\varphi\), that is, iff
\(v\entails\varphi\); and \(v\entails\varphi\lor\varphi\) iff
\(v\entails\varphi\) or \(v\entails\varphi\), that is, iff
\(v\entails\varphi\).  By (F1) all six laws hold.
\end{proof}
\end{annexitem}

\begin{annexitem}{annex:p02-w02-p02-007-bd}{thm:p02-w02-p02-007-bounds}{Bounds and complements}
\begin{proof}
\emph{Units.}  By (F3) \(v\entails\topv\), so \(v\entails\varphi\land\topv\)
iff \(v\entails\varphi\).  By (F2) \(v\not\entails\bottom\), so
\(v\entails\varphi\lor\bottom\) iff \(v\entails\varphi\).

\emph{Annihilators.}  \(v\not\entails\varphi\land\bottom\) since its right
conjunct fails, and \(v\not\entails\bottom\); both sides are false.
\(v\entails\varphi\lor\topv\) since its right disjunct holds by (F3), and
\(v\entails\topv\); both sides are true.

\emph{Complements.}  \(v\entails\varphi\land\neg\varphi\) would require
\(v\entails\varphi\) and \(v\not\entails\varphi\), which is impossible, so
the left side is false, as is \(\bottom\).  Either \(v\entails\varphi\) or
\(v\not\entails\varphi\); in the second case \(v\entails\neg\varphi\).
So \(v\entails\varphi\lor\neg\varphi\), as \(v\entails\topv\).

\emph{Negated verum.}  \(\neg\topv\) is \(\neg\neg\bottom\).  By (F3)
\(v\entails\topv\), so \(v\not\entails\neg\topv\); also
\(v\not\entails\bottom\).  Alternatively this is the instance
\(\varphi=\bottom\) of \cref{thm:p02-w02-p02-007-double-negation}.  By (F1)
every displayed law holds.
\end{proof}
\end{annexitem}

\begin{annexitem}{annex:p02-w02-p02-007-ni}{thm:p02-w02-p02-007-negated-implication}{Negated-implication law}
\begin{proof}
By the negation clause, \(v\entails\neg(\varphi\to\psi)\) iff
\(v\not\entails\varphi\to\psi\), which by (F2) holds iff
\(v\entails\varphi\) and \(v\not\entails\psi\).  By the negation clause
\(v\not\entails\psi\) iff \(v\entails\neg\psi\), so the condition is
\(v\entails\varphi\) and \(v\entails\neg\psi\), which by the conjunction
clause holds iff \(v\entails\varphi\land\neg\psi\).  By (F1) the law holds.
\end{proof}
\end{annexitem}

\begin{annexitem}{annex:p02-w02-p02-007-taut}{thm:p02-w02-p02-007-tautological-clause}{Validity of clauses}
\begin{proof}
By \cref{lem:prop-clause-satisfaction}(1), for every valuation \(v\):
\(v\entails\operatorname{form}_{\mathsf{CPL}}(C)\) iff \(v\entails L\)
for some literal \(L\in C\).

Suppose \(p\in C\) and \(\neg p\in C\).  For any \(v\), either
\(v\entails p\), or \(v\not\entails p\) and then \(v\entails\neg p\).  In
either case some literal of \(C\) holds under \(v\), so
\(v\entails\operatorname{form}_{\mathsf{CPL}}(C)\).  As \(v\) was
arbitrary, \(\entails\operatorname{form}_{\mathsf{CPL}}(C)\).

Conversely, suppose that no variable \(p\) has both \(p\in C\) and
\(\neg p\in C\).  Define a valuation \(v\) by \(v(p)=\mathbf T\) if
\(\neg p\in C\) and \(v(p)=\mathbf F\) otherwise.  Let \(L\in C\).  If \(L\)
is a variable \(p\), then \(\neg p\notin C\) by the hypothesis, so
\(v(p)=\mathbf F\) and \(v\not\entails L\).  If \(L\) is \(\neg p\), then
\(v(p)=\mathbf T\), so \(v\not\entails\neg p\).  Thus no literal of \(C\)
holds under \(v\), and \(v\not\entails\operatorname{form}_{\mathsf{CPL}}(C)\);
the formula is not valid.  The empty clause has no members, so it falls
under this case and is not valid; this agrees with
\(\operatorname{form}_{\mathsf{CPL}}(\varnothing)=\bottom\).

For the consequence, a valuation satisfies a clause set iff it satisfies
each member.  If \(C\) is tautological, every valuation satisfies \(C\), so
satisfying every member of \(S\) is equivalent to satisfying every member of
\(S\setminus\{C\}\).
\end{proof}
\end{annexitem}

\begin{annexitem}{annex:p02-w02-p02-007-cnf}{thm:p02-w02-p02-007-cnf-size}{CNF lower bound}
\begin{proof}
\emph{Step 0: the finite disjunction.}  For formulas
\(\alpha_1,\ldots,\alpha_n\) with \(n\geq1\), let \(D_n=\alpha_n\) and
\(D_i=(\alpha_i\lor D_{i+1})\) for \(1\leq i<n\), as in the convention of
Section~\ref{sec:prop-syntax-semantics-01}, so that
\(\bigvee_{i=1}^{n}\alpha_i=D_1\).  By downward induction on \(i\),
\(v\entails D_i\) iff \(v\entails\alpha_j\) for some \(j\) with
\(i\leq j\leq n\): the case \(i=n\) is immediate, and the case \(i<n\)
follows from the disjunction clause and the induction hypothesis.  With
\(\alpha_i=(P_i\land Q_i)\) and the conjunction clause,
\[
 v\entails\theta_n
 \quad\text{iff}\quad
 v(P_i)=v(Q_i)=\mathbf T\text{ for some }i\in\{1,\ldots,n\}.
 \tag{\(\ast\)}
\]
The count of connectives is direct: the formula has the \(n\) conjunctions
\(P_i\land Q_i\) and the \(n-1\) disjunctions introduced by
\(D_1,\ldots,D_{n-1}\).

\emph{Step 1: the clauses \(K_s\) are distinct, and there are \(2^n\) of
them.}  Each \(s\in S_n\) is fixed by \(n\) independent choices between
the two distinct variables \(P_i\) and \(Q_i\), so \(S_n\) has \(2^n\)
members.  Let \(s\neq s'\) and choose \(i\) with \(s(i)\neq s'(i)\).
Since the \(2n\) variables are distinct, \(s(i)\) equals \(s'(j)\) only if
\(j=i\); and \(s'(i)\neq s(i)\).  Hence \(s(i)\in K_s\) and
\(s(i)\notin K_{s'}\), so \(K_s\neq K_{s'}\).

\emph{Step 2: clause 1.}  Let \(\gamma\) be a conjunctive normal form whose
clause list \(C_1,\ldots,C_m\) enumerates the clauses \(K_s\), each once;
here \(m=2^n>0\).  By \cref{lem:prop-clause-satisfaction}(2) and (1),
\(v\entails\gamma\) iff for every \(s\in S_n\) there is \(i\) with
\(v(s(i))=\mathbf T\).  If \(v(P_{i_0})=v(Q_{i_0})=\mathbf T\) for some
\(i_0\), then for every \(s\) the variable \(s(i_0)\) is true, so
\(v\entails\gamma\).  If instead for every \(i\) one of \(P_i,Q_i\) is false
under \(v\), let \(s(i)=P_i\) when \(v(P_i)=\mathbf F\) and \(s(i)=Q_i\)
otherwise; then \(v(s(i))=\mathbf F\) for every \(i\), so no literal of
\(K_s\) holds and \(v\not\entails\gamma\).  Thus \(v\entails\gamma\) iff the
right side of \((\ast)\) holds, iff \(v\entails\theta_n\).  So
\(\gamma\equiv\theta_n\).

\emph{Step 3: clause 2, the formula is not \(\topv\).}  Let \(\gamma\) be a
conjunctive normal form with \(\gamma\equiv\theta_n\).  The valuation that
makes every variable false falsifies \(\theta_n\) by \((\ast)\), while
\(\topv\) is true under every valuation by (F3).  So \(\gamma\neq\topv\),
and by the definition of conjunctive normal form \(\gamma=N_1\) for some
clause list \(C_1,\ldots,C_m\) with \(m>0\).

\emph{Step 4: one falsified clause for each \(s\).}  For \(s\in S_n\) write
\(\bar s(i)\) for the member of \(\{P_i,Q_i\}\) other than \(s(i)\).  Define
\(v_s\) by \(v_s(s(i))=\mathbf F\) and \(v_s(\bar s(i))=\mathbf T\) for
every \(i\), and \(v_s(p)=\mathbf F\) for every other variable; this is
well defined because the \(2n\) variables are distinct.  Every
\(P_i\land Q_i\) has a false conjunct under \(v_s\), so
\(v_s\not\entails\theta_n\) by \((\ast)\), and therefore
\(v_s\not\entails\gamma\).  By \cref{lem:prop-clause-satisfaction}(2) there
is some \(j\) with
\(v_s\not\entails\operatorname{form}_{\mathsf{CPL}}(C_j)\); let \(j(s)\) be
the least such \(j\).  By \cref{lem:prop-clause-satisfaction}(1), every
literal of \(C_{j(s)}\) is false under \(v_s\).

\emph{Step 5: \(C_{j(s)}\) contains every \(s(i)\).}  Suppose
\(s(i)\notin C_{j(s)}\) for some \(i\).  Let \(w\) agree with \(v_s\) except
that \(w(s(i))=\mathbf T\).  Then \(w(P_i)=w(Q_i)=\mathbf T\), so
\(w\entails\theta_n\) by \((\ast)\), hence \(w\entails\gamma\), and by
\cref{lem:prop-clause-satisfaction}(2) and (1) some literal \(L\in C_{j(s)}\)
holds under \(w\).  That literal is false under \(v_s\).  A literal has
exactly one variable, and \(w\) and \(v_s\) differ only at \(s(i)\), so by
\cref{lem:prop-coincidence} the variable of \(L\) is \(s(i)\); thus \(L\) is
\(s(i)\) or \(\neg s(i)\).  Since \(w(s(i))=\mathbf T\), the literal
\(\neg s(i)\) is false under \(w\), so \(L=s(i)\in C_{j(s)}\), contrary to
the supposition.

\emph{Step 6: \(s\mapsto j(s)\) is injective.}  Let \(s\neq s'\) and choose
\(i\) with \(s(i)\neq s'(i)\), so that \(s'(i)=\bar s(i)\) and
\(v_s(s'(i))=\mathbf T\).  The literal \(s'(i)\) is therefore true under
\(v_s\), and since every literal of \(C_{j(s)}\) is false under \(v_s\),
\(s'(i)\notin C_{j(s)}\).  By Step~5 applied to \(s'\),
\(s'(i)\in C_{j(s')}\).  Hence \(C_{j(s)}\neq C_{j(s')}\), and in particular
\(j(s)\neq j(s')\).

The map \(s\mapsto C_{j(s)}\) is thus an injection from the \(2^n\)-element
set \(S_n\) into the set of clauses occurring in the list
\(C_1,\ldots,C_m\).  So the list contains at least \(2^n\) distinct clauses,
and \(m\geq2^n\).
\end{proof}
\end{annexitem}
